% =========================================================================
\section{Geometric work: inventory, comparison and efficiency}\label{sec:work}

Throughout this section $p\ge3$ is an integer (or the numerator $\mu$ of a rational degree,
Proposition~\ref{prop:rational}), the state $R(s)$ is supplied, and counts follow
Section~\ref{sec:convention}. For every support the power stage uses $2p-1$ parallels and the join
$L_1B$, and transport adds the join $ME$ and two parallels (the parallel through $P$ and the horizontal
readout). The supports differ in what they add: nothing for MB and AK, whose support lines are
prepared; one arc and one join for centered and decentered AD; two joins for projective AD; and three
arcs and one join for the three-mean support.

\subsection{Inventory of the three-mean support}

Fixed preparation consists of the four verticals $v,v_1,v_2,v_3$, the horizontal through $M$ and the
points $F,G_2,G_3$ of Section~\ref{sec:mean-construction}, in a rectangle of the shape
\eqref{eq:aspect}. The residual is not transferred to a rail: the line $ME$, drawn for the transport,
carries it to the three verticals at no cost.

\begin{table}[h]
\centering\small
\caption{Per-step upper inventory of the three-mean support, from $R(s)$ to $R(s_+)$, integer
$p\ge3$.}\label{tab:mean-inventory}
\begin{tabular}{@{}lccc@{}}
\toprule
Stage & $\jn$ & $\pl$ & $\ar$\\
\midrule
Triangle power endpoint $E=R(s^{p})$ & 1 & $2p-1$ & 0\\
Arc 1: centre $J_1$, opening $J_3G_3$, onto $v$ at $D_1$ & 0 & 0 & 1\\
Arc 2: centre $J_2$ through $G_2$, onto the horizontal of $M$ at $Z$ & 0 & 0 & 1\\
Arc 3: centre $D_1$, opening $FZ$, onto $v$ at $D_M$ & 0 & 0 & 1\\
Support join $AD_M$ & 1 & 0 & 0\\
Join $ME$, transport parallel, horizontal readout & 1 & 2 & 0\\
\midrule
Total & 3 & $2p+1$ & 3\\
\bottomrule
\end{tabular}
\end{table}

\begin{proposition}[Three-mean work bound]\label{prop:mean-inventory}
For integer $p\ge3$, $X>0$, a rectangle satisfying \eqref{eq:aspect} and a supplied state $R(s)$, the
protocol of Section~\ref{sec:mean-construction} realizes the three-mean correction with at most
\begin{equation}\label{eq:mean-inventory}
3\jn+(2p+1)\pl+3\ar=(4p+5)\jn+(6p+6)\ar,
\end{equation}
that is $10p+11$ traces per step; $41$ traces for cube roots. In an arbitrary rectangle one more join
suffices.
\end{proposition}
\begin{proof}
Section~\ref{sec:mean-construction} shows that the three arcs meet their lines at every state and that
$AD_M$ has the slope $\tfrac{H}{WX}[1+(p-1)q_M]$, so that \eqref{eq:support} gives $\Tm_M$. The points
$J_1,J_2,J_3$ and the two compass openings are intersections and distances of drawn or prepared
objects, which are not charged. Adding the six rows gives the left side of
\eqref{eq:mean-inventory}, and $\pl=2\jn+3\ar$ gives the right side. For the last assertion replace
$J_1$ by the intersection of $v_1'$, the vertical at distance $K_n\sqrt{\zeta^{2}-1}$ from $v$, with
the join of $E$ to a prepared centre chosen so that this intersection has the height
$H-\tfrac\lambda\gamma[1+(p-1)m(1+t)]$.
\end{proof}

The bound relies on the convention of Section~\ref{sec:convention} more than the other inventories do:
it uses seven prepared objects and three free intersections where the decentered arc uses three
prepared objects, and it prescribes the shape of the rectangle. The comparisons below should be read
with this in mind (Section~\ref{sec:limits}).

\subsection{Six supports at a common state}

\begin{table}[h]
\centering\small
\caption{Mobile work per step with the triangle power stage, integer $p\ge3$.}\label{tab:mobile}
\begin{tabular}{@{}lccccll@{}}
\toprule
Construction & Order & $\jn$ & $\pl$ & $\ar$ & Traces & Scalar scope\\
\midrule
MB rectangle & 1 & 2 & $2p+1$ & 0 & $10p+7$ & $X>1$\\
AK rectangle & 2 & 2 & $2p+1$ & 0 & $10p+7$ & global\\
AD rectangle & 3 & 3 & $2p+1$ & 1 & $10p+9$ & global\\
Projective AD & 4 & 4 & $2p+1$ & 0 & $10p+9$ & global\\
Decentered AD & 5 & 3 & $2p+1$ & 1 & $10p+9$ & global, $p>2$\\
Three means & 9 & 3 & $2p+1$ & 3 & $10p+11$ & global, $p>2$\\
\bottomrule
\end{tabular}
\end{table}

For cube roots the per-step bounds are $37,37,39,39,39,41$. Table~\ref{tab:cubic} runs all six from the
same state $X=2$, $s_0=3/4$ with the same stopping test on the reciprocal relative error.

\begin{table}[h]
\centering\small
\caption{Steps / summed trace bounds to reach $|s/\alpha-1|\le10^{-d}$ for $p=3$, $X=2$, $s_0=3/4$.
Each total is the per-step bound times the number of steps. Counts agree at 250 and 500 working
digits.}\label{tab:cubic}
\begin{tabular}{@{}lccccc@{}}
\toprule
Construction & Order & $d=6$ & $d=12$ & $d=30$ & $d=60$\\
\midrule
MB rectangle & 1 & 8 / 296 & 17 / 629 & 44 / 1628 & 90 / 3330\\
AK rectangle & 2 & 3 / 111 & 4 / 148 & 5 / 185 & 6 / 222\\
AD rectangle & 3 & 2 / 78 & 3 / 117 & 3 / 117 & 4 / 156\\
Projective AD & 4 & 2 / 78 & 2 / 78 & 3 / 117 & 3 / 117\\
Decentered AD & 5 & 1 / 39 & 2 / 78 & 2 / 78 & 3 / 117\\
Three means & 9 & 1 / 41 & 1 / 41 & 2 / 82 & 2 / 82\\
\bottomrule
\end{tabular}
\end{table}

Higher order alone does not settle the comparison at a fixed accuracy. For cube roots the three-mean
support takes $1,1,2,2$ steps against $1,2,2,3$ for the decentered arc, that is $41,41,82,82$ traces
against $39,78,78,117$: it wins at $d=12$ and $d=60$ and loses by two traces at $d=6$ and $d=30$, where
both need the same number of steps.

\subsection{Efficiency index}

For an iteration of order $\varrho$ costing $N$ units per step, the classical efficiency index is
$\varrho^{1/N}$~\cite{ostrowski,traub64}: the number of correct digits is asymptotically multiplied by
$\varrho^{1/N}$ per unit of work. We compare its logarithm $\iota=\log\varrho/N(p)$, with $N(p)$ the
trace bound of Table~\ref{tab:mobile}.

\begin{proposition}[Efficiency ordering]\label{prop:index}
For every integer $p\ge3$,
$\iota_{\mathrm{AK}}<\iota_{\mathrm{AD}}<\iota_{\mathrm{proj}}<\iota_{\mathrm{dec}}<\iota_{M}$, and the
last inequality holds for every real $p>0$. The three-mean support has the largest index of the six
supports at every integer degree $p\ge3$.
\end{proposition}
\begin{proof}
The last three of the first group share $N=10p+9$ and have increasing orders, and
$\log2/(10p+7)<\log3/(10p+9)$ because $10p\log(3/2)>9\log2-7\log3$, whose right side is negative.
Finally $\log9/(10p+11)>\log5/(10p+9)$ reads $10p\log(9/5)>11\log5-9\log9$, whose right side is
negative because $5^{11}<9^{9}$. MB has order one and index zero.
\end{proof}

\begin{table}[h]
\centering\small
\caption{$10^{3}\iota=10^{3}\log\varrho/N(p)$. The largest value in each column is in bold.}\label{tab:index}
\begin{tabular}{@{}lccccccc@{}}
\toprule
Construction & $p=3$ & $5$ & $8$ & $12$ & $13$ & $20$ & $50$\\
\midrule
AK & 18.73 & 12.16 & 7.97 & 5.46 & 5.06 & 3.35 & 1.37\\
AD & 28.17 & 18.62 & 12.34 & 8.52 & 7.90 & 5.26 & 2.16\\
Projective AD & 35.55 & 23.50 & 15.58 & 10.75 & 9.97 & 6.63 & 2.72\\
Decentered AD & 41.27 & 27.28 & 18.08 & 12.48 & 11.58 & 7.70 & 3.16\\
Three means & \textbf{53.59} & \textbf{36.02} & \textbf{24.15} & \textbf{16.77} & \textbf{15.58} & \textbf{10.41} & \textbf{4.30}\\
\bottomrule
\end{tabular}
\end{table}

The reason is visible in Table~\ref{tab:mobile}: the power stage costs about $10p$ traces for every
support, and the three-mean support adds only two traces to the decentered arc, so it leads at every
degree. The ratio of its index to that of the decentered arc is $1.30$ at $p=3$ and tends to
$\log9/\log5\approx1.37$. The index is an asymptotic
quantity built from upper inventories; at a finite accuracy the integer number of steps makes the
comparison lumpy, as Table~\ref{tab:p20} shows.

\begin{table}[h]
\centering\small
\caption{Steps / summed trace bounds for $p=20$, $X=2$, $s_0=3/4$ (per-step bounds $209$ and
$211$).}\label{tab:p20}
\begin{tabular}{@{}lccccc@{}}
\toprule
Construction & $d=30$ & $d=60$ & $d=250$ & $d=1000$ & $d=2500$\\
\midrule
Decentered AD & 4 / 836 & 5 / 1045 & 5 / 1045 & 6 / 1254 & 7 / 1463\\
Three means & 3 / 633 & 3 / 633 & 4 / 844 & 4 / 844 & 5 / 1055\\
\bottomrule
\end{tabular}
\end{table}

At $p=20$ the three-mean support saves one or two steps at every target of the table, and has the
smaller total at each of them.

% =========================================================================
\section{Direct root enclosures for real degrees}\label{sec:enclosure}

The power stage can be avoided altogether by constructing the root directly. The input is a positive
length $X$ and a real degree $p$, also given as a length; the output is an approximation of
$r=X^{1/p}$ with prescribed relative error. A real input length is assumed to be available: this does
not assert that every real number is constructible from a unit segment. This calculator is the only one in the paper whose every
step is a finite construction for irrational $p$. We first treat $X>1$ and $p\ge1$; the cases $X=1$ and
$p=1$ are immediate, and Proposition~\ref{prop:allp} removes both restrictions.

\subsection{One mean proportional per subdivision}

Set $\theta=1/p$, $a=1$, $b=X$ and $\lambda=\theta$. Here $a,b$ are moving rail values. The invariant is
\begin{equation}\label{eq:inv}
0<a<b,\qquad 0\le\lambda\le1,\qquad r=a^{1-\lambda}b^{\lambda}.
\end{equation}
If $\lambda\in\{0,1\}$ return the corresponding endpoint, which is exact. Otherwise construct
$m=\sqrt{ab}$ with the classical mean-proportional semicircle~\cite{euclid} and replace
\begin{equation}\label{eq:subdiv}
(a,b,\lambda)\longmapsto
\begin{cases}(a,\,m,\,2\lambda),&\lambda\le\tfrac12,\\[2pt](m,\,b,\,2\lambda-1),&\lambda>\tfrac12.\end{cases}
\end{equation}
Both branches preserve \eqref{eq:inv}. Thus $p$ determines the choices and the weight, never the length
of a power chain. After $n$ subdivisions
\begin{equation}\label{eq:gap}
\frac ba=X^{2^{-n}},\qquad z:=\frac ba-1\ge0.
\end{equation}
No logarithm, exponential or power $s^{p}$ is evaluated. The branch test compares two constructed
lengths, $\lambda$ and $\tfrac12$; for irrational $p$ equality never occurs, and for dyadic $1/p$ the
process stops with an exact endpoint.

\begin{figure}[h]
\centering\figSubdivision
\caption{A subdivision constructs one mean proportional. The semicircle has diameter $OS=a+b$, centre
$C$ and junction $OT=a$; the orange perpendicular $TG$ is the mean $m$. The same device is rotated onto the vertical rail in the counted protocol.
The figure uses $a=1.2$, $b=2.8$ only to show the incidence.}\label{fig:subdiv}
\end{figure}

\subsection{A fifth-order final enclosure}

The harmonic--arithmetic readout of the companion gallery~\cite{gallery},
$L_2=ab/[(1-\lambda)b+\lambda a]$ and $U_2=(1-\lambda)a+\lambda b$, encloses $r$ by the weighted
arithmetic--geometric mean inequality, with
\[ \frac{U_2}{L_2}-1=\lambda(1-\lambda)\frac{z^{2}}{1+z}\le\sinh^{2}\!\Bigl(\frac{\log X}{2^{n+1}}\Bigr), \]
a width quadratic in $z$. A single rational readout of degree two gives a fifth-order width. It is the
diagonal $[2/2]$ Pad\'e approximant of $(1+z)^{\tau}$ at zero~\cite{dlmf,baker}. For $0\le\tau\le1$ put
\begin{equation}\label{eq:pade}
\Pi_\tau(z)=12+6(2+\tau)z+(\tau+1)(\tau+2)z^{2},\qquad
\psi_\tau(z)=\frac{\Pi_\tau(z)}{\Pi_{-\tau}(z)},
\end{equation}
whose denominators are positive for $z\ge0$, and
\begin{equation}\label{eq:L5}
L_5=a\,\psi_\lambda(z),\qquad U_5=\frac{b}{\psi_{1-\lambda}(z)}.
\end{equation}

\begin{proposition}[Enclosure and finite termination]\label{prop:enclosure}
For $X\ge1$, real $p\ge1$ and every subdivision depth,
\begin{equation}\label{eq:width}
L_5\le X^{1/p}\le U_5,\qquad
\frac{U_5}{L_5}-1=\frac{\lambda(1-\lambda)(2-\lambda)(1+\lambda)\,z^{5}}{\Pi_\lambda(z)\,\Pi_{1-\lambda}(z)}
\le\frac{z^{5}}{256}.
\end{equation}
Given $\eps>0$, choose a positive rational $\beta$ with $\beta^{5}/256\le\eps$, stop when
$b\le(1+\beta)a$ and return $L_5$. The procedure terminates and its relative error is at most $\eps$.
A sufficient depth depends on $X$ and $\eps$ only, not on $p$.
\end{proposition}
\begin{proof}
For $0\le\tau\le1$, direct differentiation gives
\[ (1+z)\bigl(\Pi_\tau'\Pi_{-\tau}-\Pi_\tau\Pi_{-\tau}'\bigr)-\tau\,\Pi_\tau\Pi_{-\tau}
=-\tau(1-\tau^{2})(4-\tau^{2})\,z^{4}\le0. \]
Hence $\log\psi_\tau(z)-\tau\log(1+z)$ is nonincreasing from zero, so
$\psi_\tau(z)\le(1+z)^{\tau}$. With $\tau=\lambda$ this is $L_5\le r$; with $\tau=1-\lambda$ and
$(1+z)^{\lambda}(1+z)^{1-\lambda}=1+z$ it is $r\le U_5$. Multiplying out gives the equality in
\eqref{eq:width}. With $\chi=\lambda(1-\lambda)\le\tfrac14$ the numerator coefficient is
$\chi(2+\chi)\le\tfrac9{16}$, and $\Pi_\lambda,\Pi_{1-\lambda}\ge12$, so the width is at most
$z^{5}/256$. At the stopping test $0\le z\le\beta$, so $r/L_5-1\le\eps$ and $(r-L_5)/r\le\eps/(1+\eps)$. By \eqref{eq:gap}, $z\to0$,
so the test is eventually met; any $n$ with $2^{n}\log(1+\beta)\ge\log X$ suffices. Equivalently
$X\le(1+\beta)^{2^{n}}$, which can be checked in exact rational arithmetic for rational input.
\end{proof}

It suffices to draw $L_5$ alone; the analytic upper bound certifies its accuracy without constructing
a second Pad\'e quotient. For $0<X<1$, constructing and inverting this lower bound for $1/X$ gives an
upper approximation of the original root with relative error at most $\eps$. The width is of order
five in the rail gap $z$; this is not a fifth-order iteration, since each subdivision only halves the
logarithmic gap.

\begin{proposition}[Every nonzero real degree and every positive input]\label{prop:allp}
Let $p\neq0$ be real and $X>0$, $X\neq1$.
\begin{enumerate}[label=(\roman*)]
\item If $X>1$ and $0<p<1$, let $k=\lfloor1/p\rfloor\ge1$ and start from $a=X^{k}$, $b=X^{k+1}$,
  $\lambda=1/p-k$. Then \eqref{eq:inv}, \eqref{eq:gap} and Proposition~\ref{prop:enclosure} hold
  unchanged.
\item If $0<X<1$, apply the construction to $1/X$ and invert: the enclosure $[1/U_5,1/L_5]$ has the same
  relative width.
\item If $p<0$, apply the construction to $|p|$ and invert the output.
\end{enumerate}
\end{proposition}
\begin{proof}
In (i), $a^{1-\lambda}b^{\lambda}=X^{k+\lambda}=X^{1/p}$, $0\le\lambda<1$ and $b/a=X$, so the invariant
and the gap formula are those of the case $p\ge1$, and the proof of Proposition~\ref{prop:enclosure}
uses nothing else. Inversion maps an enclosure $[L,U]$ to $[1/U,1/L]$ with the same ratio.
\end{proof}

\subsection{A complete, conservative drawing inventory}\label{sec:macros}

Use the extended vertical rail $R(q)$, its origin $R(0)$ and the prepared centres of
Section~\ref{sec:rect}; $X$ and $p$ are given on the rail. The macros are:
\begin{itemize}
\item \emph{Sum or nonnegative difference} of two rail values: one compass arc on the rail.
\item \emph{Product or quotient}: five joins (Section~\ref{sec:rect}); by a prepared constant, two joins.
\item \emph{Mean proportional}, $2\jn+\pl+3\ar$: form $R(a+b)$ with one arc and $R((a+b)/2)$ with the
  two-join fixed scale; draw the semicircle through $R(0)$; draw the horizontal through $R(a)$ with one
  parallel macro; transfer the resulting length to the rail with one arc.
\end{itemize}
The updated weight uses the fixed scale $2\lambda$ and, on a right branch, one subtraction. A
subdivision with its weight update costs at most $4\jn+\pl+4\ar$, and a stopping test, which scales $a$
by the prepared $1+\beta$, two joins. For the final readout use
\begin{equation}\label{eq:L5eval}
L_5=a\Bigl(1+\lambda\,\frac{6z(2+z)}{\Pi_{-\lambda}(z)}\Bigr),\qquad
\Pi_{-\lambda}(z)=12+6(2-\lambda)z+(1-\lambda)(2-\lambda)z^{2}.
\end{equation}
One evaluation schedule uses seven general products, namely $(1-\lambda)(2-\lambda)$, $z^{2}$,
$(2-\lambda)z$, $(1-\lambda)(2-\lambda)z^{2}$, $z(2+z)$, the weighted quotient and the final scaling by
$a$; two divisions, including $b/a$; two fixed scalings by $6$; and seven sums or differences. Its cost
is at most $49\jn+7\ar$.

\begin{table}[h]
\centering\small
\caption{Trace inventory of the direct enclosure for $X>1$, $p\ge1$.}\label{tab:subdiv}
\begin{tabular}{@{}lccc@{}}
\toprule
Operation & $\jn$ & $\pl$ & $\ar$\\
\midrule
Initial weight $\lambda=1/p$ & 5 & 0 & 0\\
$n$ subdivisions and weight updates & $4n$ & $n$ & $4n$\\
At most $n+1$ stopping tests & $2n+2$ & 0 & 0\\
One final lower Pad\'e readout & 49 & 0 & 7\\
\midrule
Total, parallels expanded & $8n+56$ & 0 & $7n+7$\\
\bottomrule
\end{tabular}
\end{table}

Thus $15n+63$ traces suffice for $X>1$, $p\ge1$. The other cases of Proposition~\ref{prop:allp} add:
$5k+1$ traces in case (i), for the $k$ products $X^{2},\dots,X^{k+1}$ and the difference $1/p-k$;
two five-join inversions in case (ii), giving $15n+73$ for $0<X<1$, $p\ge1$; one five-join
inversion in case (iii). Finding $k$ is a decision and is not charged. Constructing both endpoints
$L_5$ and $U_5$ would require additional work and is not included in this single-output count.

\subsection{Depth, precision and comparison with the iterations}

\begin{table}[h]
\centering\small
\caption{Sufficient depths at $X=2$, uniformly in real $p\ge1$. The Pad\'e depths use
$\beta=19/100,\ 3/40,\ 3/250,\ 3/1000$; the inequalities $\beta^{5}/256\le\eps$ and
$(1+\beta)^{2^{n}}\ge2$ verify them in rational arithmetic.}\label{tab:depth}
\begin{tabular}{@{}cccc@{}}
\toprule
Relative width target & Harmonic--arithmetic depth & Pad\'e depth & Trace bound $15n+63$\\
\midrule
$10^{-6}$ & 9 & 2 & 93\\
$10^{-8}$ & 12 & 4 & 123\\
$10^{-12}$ & 19 & 6 & 153\\
$10^{-15}$ & 24 & 8 & 183\\
\bottomrule
\end{tabular}
\end{table}

\begin{figure}[h]
\centering\figDepth
\caption{Sufficient worst-case depths of Table~\ref{tab:depth}, uniformly in $p\ge1$. Points are
verified depth certificates; the connecting segments only guide the eye and do not represent fractional
subdivisions.}\label{fig:depth}
\end{figure}

For $p=37.5$, four subdivisions yield $a=1$, $b=2^{1/16}$ and $\lambda=32/75$, and the root lies between
$1.0186558096$ and $1.0186558103$; these rounded endpoints can be checked exactly by raising them to the
power $75$ and comparing with $4$. At $X=2$ and $\eps=10^{-8}$ the construction has the bound
$88\jn+35\ar=123$ traces.

\paragraph{When to use which calculator.}
The subdivision gains a fixed number of bits per step, the iterations multiply the number of correct
digits. For cube roots of $2$, the depths needed for width $10^{-d}$ with $d=6,12,30,60$ are
$n=2,6,18,38$, that is $93,153,333,633$ traces, against $39,78,78,117$ for the decentered arc and
$41,41,82,82$ for the three-mean support in Table~\ref{tab:cubic}. (The two columns are not the same quantity: a guaranteed width on one side, an
observed error on the other.) For small integer or rational degrees and high accuracy the iterations
are cheaper; for irrational degrees, and for rational degrees with a large numerator, the subdivision
is the calculator to use (Remark~\ref{rem:pi}).
